\chapter{Related work}
\label{ch:lit}

% Written 21 September 2026 from analysis/SPINE.md section 2.
% TODO(lit): two citations are wanted and are not yet in refs.bib. Section 2.1 should cite the
% two-timescale / Stackelberg learning-dynamics literature (Fiez, Chasnov and Ratliff 2020; Borkar
% 1997) and Section 2.4 the commitment argument in Chicken (Schelling 1960). Add them only after the
% usual verification; the text below stands without them.

This chapter places the thesis in three literatures: opponent shaping, from
differentiable methods to prompt-based ones; language-model agents in social
dilemmas and commons; and the resource economics the environment is drawn
from.
The gap it identifies is not a combination of attributes that no one has
tried.
It is an identification problem: published evidence that a shaping objective
changed a co-player's learning comes from comparisons in which the shaper
also learns on a different timescale, and in the games used, that difference
predicts the same outcome.

\section{Opponent shaping}
\label{sec:lit-shaping}

Foerster et al.\ \cite{foerster2018} introduce opponent shaping: an agent
differentiates through its co-player's anticipated learning step.
The method is white-box and one-step.
Language-model co-players have neither accessible parameters nor a known
learning rule, so that channel is closed.

Lu et al.\ \cite{lu2022} recast shaping as model-free reinforcement learning
in a meta-game: a meta-step is an inner episode, the meta-action is the
shaper's next inner policy, and the meta-reward is the inner return.
Khan et al.\ \cite{khan2024} scale the same construction to long-horizon
gridworlds and separate two kinds of memory, history within an episode and
context across episodes, finding both necessary.
They also warn that in CoinGame the current state leaks past actions, so an
agent can appear to shape without conditioning on anything across episodes.
A resource stock is itself a statistic of past harvests, so the same warning
applies here, and it is why one arm removes the previous-round line and
another removes the trial-memory prompt (Chapter~\ref{ch:design}).

Both of those meta-game methods evaluate a shaper by freezing it and
training a fresh co-player against it.
That test is adopted here as E1 (Chapter~\ref{ch:method}); it is not part of
the evaluation in the language-model work below, which replays the
co-trained pair.

Garc\'ia Segura et al.\ \cite{garciasegura2025} adapt model-free opponent
shaping to language-model agents.
A shaper and a naive learner, both Gemma-2-2b-it with LoRA adapters, play
five repeated \(2\times2\) games through prompts; the shaper updates once per
trial on the trial's return while the naive learner updates after every
episode.
The reported evidence is average reward per step and the visitation of
joint-action cells, and in the iterated Chicken game the shaper reaches an
exploitative asymmetry.
Two features of that comparison matter here.
The shaper's learning rate is an order of magnitude below its co-player's
and it takes one update where the co-player takes five, so the two agents
differ in three respects at once; and the paper's shaping control enriches
the naive learner's observations rather than matching the shaper's schedule.
In games where commitment pays, a slower learner is closer to a committed
player, so the timescale difference alone predicts an exploitative
asymmetry.
Distinguishing the two explanations requires a control that holds learning
rate, clip, batch and update schedule fixed and removes only the
cross-episode objective; that control is the trial-batched arm of
Chapter~\ref{ch:design}, and it is the methodological contribution of this
thesis.
A second feature matters for any transfer of the method to episodic games:
the released estimator chains generalised advantage estimation across the
episodes of a trial through one sequence per parallel game, which is sound
where a repeated matrix game has no terminal state, and which
Chapter~\ref{ch:results} shows bootstraps through the environment's reset
where an episode can end by collapse.

Duque et al.\ \cite{duque2025} collapse shaping into a policy-gradient term
that multiplies own and opponent advantages, and evaluate it on domains that
include a spatial commons with probabilistic regrowth; a later application
\cite{duque2026} places the same method in a climate-investment simulator.
Both assume observable opponent advantages, which the prompt-only regime
does not provide, and both train in symmetric self-play rather than against
an independent learner.
Piche et al.\ \cite{piche2025} train language-model agents with advantage
alignment in self-play on social dilemmas, and use group-relative common
random numbers to reduce variance; the paired noise tables used for
measurement here are that device, not that training method.

\section{Language-model agents in dilemmas and commons}
\label{sec:lit-llm}

Akata et al.\ \cite{akata2025} characterise frozen language models across
structurally distinct \(2\times2\) games: no agent learns, so no agent
shapes.
Tennant et al.\ \cite{tennant2025} train Gemma-2-2b-it with intrinsic moral
rewards on the iterated prisoner's dilemma, which modifies the trained
agent's own objective rather than its co-player's learning.

Piatti et al.\ \cite{piatti2024} place societies of frozen language models
on a shared stock that regrows by doubling, and find that most models
collapse it; communication and a universalisation prompt are the
interventions that prevent collapse.
No agent is fine-tuned, and the regeneration law is not the logistic one
used here.
That study is the closest existing statement about language models and a
commons, and it concerns prompting a frozen model rather than training one
against a co-learner.

P\'erolat et al.\ \cite{perolat2017} study a spatial commons with twelve
independent learners, where sustainability arrives through exclusion:
tagging and territory.
The environment here removes exclusion deliberately so that the harvest
decision is the only decision and the game can be solved exactly.
Their sustainability, efficiency and equality metrics are reported in
Chapter~\ref{ch:results} as outcomes, not as evidence about the objective;
Chapter~\ref{ch:method} explains why they cannot discriminate shaping in
this game.
Leibo et al.\ \cite{leibo2017} give the empirical-payoff-matrix reduction
that Chapter~\ref{ch:env} applies to constant strategies.

\section{Resource dynamics}
\label{sec:lit-cpr}

The regrowth law is the logistic surplus-production model of Schaefer
\cite{schaefer1954}, the standard single-stock model of fisheries economics
\cite{clark1990}; Gordon \cite{gordon1954} supplies the open-access result
that unmanaged effort dissipates rent.
Reed \cite{reed1979} studies the same stock under multiplicative
recruitment uncertainty and shows that the optimal harvest is a
constant-escapement feedback rule, and Sethi et al.\ \cite{sethi2005}
separate growth, measurement and implementation uncertainty and show that
the result survives the first alone.
The shock used here is growth uncertainty in that sense: a bounded,
mean-one, three-point multiplier on the growth increment, with the stock and
both requests observed exactly.
The parameters are design choices fixed by the invariants of
Chapter~\ref{ch:env}, not calibrations from any of these papers.

Under constant strategies the game is chicken rather than a prisoner's
dilemma, the class Rapoport and Chammah \cite{rapoport1966} name and the
class ShapeLLM already evaluates; Maynard Smith and Price
\cite{maynardsmith1973} give its evolutionary ancestor.
The relevant property here is not the name but the consequence: in chicken
the player who commits wins, which is why the \(m=3\) regime is a negative
control and why a timescale control is necessary at \(m=2\).

\section{The gap}
\label{sec:lit-gap}

Observational work on language models and commons freezes every agent
\cite{piatti2024,akata2025}.
Interventional work that trains a language-model shaper against an
independent language-model learner \cite{garciasegura2025} compares it with
learners that differ from it in learning rate and update frequency as well
as in objective, in games where that difference alone predicts the reported
outcome, and evaluates by replaying the co-trained pair rather than by
freezing the shaper.
No published design separates the objective from the timescale that
accompanies it, and none measures what the co-player acquired independently
of the partner that trained it.
This thesis supplies both, in a sequential commons whose exact solution
fixes what a co-player's restraint is worth and whose two regimes differ in
one parameter.
